\section{Robustness, mechanisms, and limits}\label{sec:robustness}

The central theorem is not a claim that the precise threshold of two agenda holders arises in every model of international standardization. Its force comes from the conjunction of two properties established in the preceding sections. First, a government cannot change the content of the public standard when exercising an agenda right: it can only bring a legally specified interface to a national vote. Second, firms understand this possibility before deciding how much irreversible compatibility capital to build. Our comparative statics hold across an explicit rectangle of the original parameters, but the public reference and political institution remain economically substantive assumptions.

\subsection{A local neighborhood of primitive parameters}\label{subsec:local}

One concern with a theory certified over a specific parameter rectangle is that its equilibrium classification could disappear under an arbitrarily small perturbation of the primitives. To examine that concern, we retain the entire original game and let all seven parameters vary in a nondegenerate product neighborhood:
\begin{equation}
\begin{gathered}
 \eta\in[.10,.14],\qquad K\in[.54,.56],\qquad
 \phi\in[.13,.30],\\
 \rho\in[.319,.321],\qquad
 \kappa\in[.019,.021],\\
 v\in[.009999,.010001],\qquad
 \gamma_0\in[.139999,.140001].
\end{gathered}
\label{eq:robust-box}
\end{equation}
For every point in this box, strict compliance, initial \(NNN\)-to-\(M\) existence, and exclusion of a privately retained pure equilibrium under \(r\le2\) and \(|\calA|\ge2\) remain valid. The first five parameters admit appreciable simultaneous perturbations relative to the last two; the intervals for \(v,\gamma_0\) are intentionally conservative. Their endpoints are not inferred economic bifurcation thresholds. They merely yield a transparent exact certificate of a genuine seven-dimensional local stability result. The constructive \emph{necessity} direction of Proposition~\ref{prop:agenda}, under which exactly one proposer can sustain the complementary bilateral RS, has not separately been re-certified over this entire seven-dimensional box.

\begin{proposition}[Local sufficiency of the agenda threshold]
\label{prop:robust}
For all primitives in (\ref{eq:robust-box}), if \(r\in\{1,2\}\)
and \(|\calA|\ge2\), every pure-strategy SPE terminal standard type
is mandatory public IS. At least one such SPE exists.
\end{proposition}

\noindent The certificate behind Proposition~\ref{prop:robust}
uses the same four corporate investment-count cases as
Proposition~\ref{prop:agenda}.
The two-developer proof requires strict signs for the
low-\(a\) national approval and high-\(a\) private derivative
polynomials. The three-developer proof requires nonnegative
policy-gain numerators in the aligned and crossed
technical-order cases. Each is represented on a unit cube
by a polynomial Bernstein basis; nonnegative coefficients
then certify the sign throughout the corresponding
continuous region. Exact-rational code verifies the
five-parameter rectangles in
\(\eta,\phi,\rho,\kappa\), together with the remaining
technology coordinates, at the endpoints of \(\phi\).
Because each sign numerator is affine in \(\phi\), endpoint
certification implies the intervening continuum.

Adding \(v\) and \(\gamma_0\) introduces a complication:
some three-developer Bernstein coefficients are exactly zero
at the boundary where all relevant private technologies
coincide with \(O(1,0)\).
Strictly positive coefficients are stable under small perturbations,
but zero coefficients require a separate argument.
The relevant boundary welfare numerators vanish identically
for \emph{all} \(v,\gamma_0\), not merely at their baseline values.
After clearing the strictly positive public-technology denominator,
every other coefficient has a rational lower bound above \(1/50\).
An explicit power-basis coefficient-norm estimate bounds
its joint \(v,\gamma_0\) perturbation below \(0.013\)
on the stated narrow intervals.
Hence none of the nonzero coefficients changes sign,
and all structural zeros remain zeros.
Appendix~\ref{app:certificates} states the algebraic objects,
normalizations, and scripts corresponding to these assertions.
The appendix is part of the proof, not an invitation to
accept a finite numerical parameter grid as sufficient.

\subsection{What carries the result}\label{subsec:mechanism}

Three features do the economic work.

\paragraph{Independent public agenda rights.}
With two authorized proposers, a two-country coalition
that gains from the mandate cannot be silenced by
a single excluded national government.
If only one government controls the agenda and that
government is outside the privately compatible pair,
its refusal blocks an otherwise successful vote.
The institutional comparison therefore changes the
set of effectively available policies, not a firm's
underlying marginal costs.

\paragraph{Continuous irreversible private design.}
An installed private technology may induce national
opposition to the public mandate. Yet it is not enough
for national opposition to exist at an arbitrary
technology profile. The firm may have paid for
compatibility or commonality that does not enter any
effective pair minimum, or may be able to trim a
costly commonality attribute without changing the
policy continuation. These private deviations sharply
restrict the histories that can be supported in a
pure initial-stage equilibrium. Without this
investment argument, the set-theoretic agenda
observation would be an incomplete equilibrium
theorem.

\paragraph{A credible public fallback.}
The public \(O(1,0)\) reference is already legally
available, and an open developer can complete its
missing compatibility stock without paying a second
fixed setup. These conditions make post-mandate
compliance strictly attractive and make the
government's policy comparison well-defined
at every private investment history. If the public
interface were patented, costly to negotiate, subject
to a design choice by governments, or usable only
through private licensing, a new continuation game
would be required. Our proof does not extend to such
systems by relabeling the current parameter values.

\subsection{Relation to existing research}\label{subsec:literature}

The literature already provides several reasons to
expect private and national standard-setting
incentives to diverge. \citet{KatzShapiro1985}
study network externalities and compatibility in
imperfect competition, and \citet{FarrellSaloner1988}
compare coordination through standard-setting
committees and markets. \citet{AuriolBenaim2000}
analyze decentralized standardization dynamics.
Thus neither network-mediated coordination problems
nor inefficient decentralization is a novel claim.

The international policy setting is particularly
close to \citet{GandalShy2001}, who study
government recognition of foreign standards and
three-country standardization unions with an excluded
nonmember. \citet{Klimenko2009} analyzes
strategic international compatibility policies
and the possibility of excessive or insufficient
compatibility, while \citet{Costinot2008} compares
mutual recognition with national-treatment rules
and finds institutional welfare rankings sensitive
to the economic environment. Our objective is
not to relabel these policy distortions as new.
The present rule variable is the set of national
governments authorized to \emph{introduce} a single
common public mandate and the quorum needed to
adopt it, after privately selected sunk
technical investments.

There is also close conceptual overlap with
\citet{FarrellSimcoe2012}, who study a
standards-setting war of attrition with neutral
intervention, screening, and ex-ante incentives
to improve proposals. Their analysis explicitly
contemplates a two-vote rule in a different
standards-governance environment.
\citet{Ruiz2004} examines governments'
foreign-standard recognition alongside
endogenous firms' product characteristics.
We therefore do not claim to be the first
to connect standard-setting rules with
innovation or to endogenize product differentiation
under government standardization policy.
The more limited proposition tested here concerns
a pure-equilibrium \emph{terminal-outcome set}
under a specific three-country national
agenda/quorum constitution and continuous
sunk open-technology stocks.

This literature comparison remains a manuscript
positioning argument, not an exhaustive
priority certification. In particular,
the full final published propositions and
proof appendices of all the closest
international-standards papers have not yet
been checked independently.
A submission decision should preserve
that qualification and complete the
remaining source audit.

\subsection{Scope of equilibrium and welfare claims}\label{subsec:limits}

The model excludes proprietary X/Y interfaces,
private recognition and licensing mechanisms,
budget-balanced intergovernmental transfers,
and endogenous public standard design.
The national tie rule selects YES and
the constitution fixes which governments
can propose. Replacing any of these may
alter continuation equilibria and invalidate
the institutional result. In addition,
the paper classifies only pure-strategy
SPE terminal standard types; it neither
enumerates all pure initial strategies
under the public equilibrium nor
rules out mixed SPEs. Finally,
the ex-ante welfare comparison ranks
two \emph{specified feasible equilibrium
investment paths}; it is not the
unconstrained first-best technology
allocation. These restrictions delimit
the intended policy interpretation.
